\section{Gowers-Hatami for BLR linearity test over $\FF_q$}

Assume $q = p^s$ be a prime with power $s\in \mathbb{Z}^+$, $n \ge 1$, and we set
\[
X := \F_q^n, \qquad A := \F_q^\times, \qquad \omega := e^{2\pi i/p}.
\]

% \begin{definition}[BLR verifier and distance]\label{def:blr}\lean{BlrPcp.BLRAcceptsAt, BlrPcp.distanceToLinear, BlrPcp.acceptanceProbabilityBLR}\leanok
% Given $f : X \to \F_q$, the verifier $\BLR^f$ samples $a, b \in A$ and
% $x, y \in X$ uniformly and accepts iff $af(x) + bf(y) = f(ax + by)$. Let the test rejection probability be bounded by
% \[
% \Pr_{a,b,x,y}\!\bigl[af(x) + bf(y) \ne f(ax + by)\bigr] \leq \varepsilon.
% \]
% For $\alpha \in X$, write $\ell_\alpha(x) := \inner{\alpha}{x}$, $\delta_\alpha := \Pr_x[f(x) \ne \ell_\alpha(x)]$,
% and $\delta(f, \lin) := \min_{\alpha} \delta_\alpha$.
% \end{definition}
% At the end of the section, we prove the BLR soundness using Gowers-Hatami theorem.

First notice that linearity implies two properties: 
\begin{enumerate}
    \item Additivity $f(x + y) = f(x) + f(y), \quad \forall x, y \in X$. 
    
    Note that for $q$ prime additivity is actually equivalent to linearity. However this is not the case for prime powers $q=p^s$ for $s > 1$.
    
    \item Homogeneity $f(a x) = a f(x), \quad \forall x \in X, \, a \in A$.

    For $q=2$, the two properties are redundant. 
    For $q > 2$, turns out testing homogeneity is crucial for a test to have a high soundness.
\end{enumerate}

% In the rest of the section, we discuss an encoding of both properties through a group construction by ChatGPT.


\subsection{The twisted group \texorpdfstring{$G$}{G}}\label{sec:group}

Define a group $G$ by the set $X \times A$ with the binary operation
\begin{equation}\label{eq:star}
(x, a) \star (y, b) := \bigl(b^{-1}x + a^{-1}y,\; ab\bigr).
\end{equation}

\begin{lemma}[$G$ is an abelian group]\label{lem:group}\lean{BlrPcp.Twist.instCommGroup}\leanok
The operation $\star$ in~\eqref{eq:star} makes $G$ into an abelian group with identity
$e = (0, 1)$ and inverses $(x, a)^{-1} = (-a^2 x, a^{-1})$. 
\end{lemma}

\begin{proof}
\emph{Associativity.} For $(x, a), (y, b), (z, c) \in G$,
\begin{align*}
\bigl((x, a) \star (y, b)\bigr) \star (z, c)
&= (b^{-1}x + a^{-1}y,\; ab) \star (z, c) = \bigl(b^{-1}c^{-1}x + a^{-1}c^{-1}y + a^{-1}b^{-1}z,\; abc\bigr),
% &= \bigl(c^{-1}(b^{-1}x + a^{-1}y) + (ab)^{-1}z,\; abc\bigr) \\
% &= \bigl(b^{-1}c^{-1}x + a^{-1}c^{-1}y + a^{-1}b^{-1}z,\; abc\bigr),
\end{align*}
and an identical computation gives the same expression for $(x, a) \star \bigl((y, b) \star (z, c)\bigr)$.

\emph{Commutativity.}
$(y, b) \star (x, a) = (a^{-1}y + b^{-1}x,\; ba) = (b^{-1}x + a^{-1}y,\; ab) = (x, a) \star (y, b).$

\emph{Identity.} For any $(y, b) \in G$,
$(0, 1) \star (y, b) = (b^{-1}\cdot 0 + 1^{-1}\cdot y,\; 1\cdot b) = (y, b)$.

\emph{Inverse.} For any $(x, a) \in G$,
\[
(x, a) \star (-a^2 x, a^{-1}) = \bigl((a^{-1})^{-1}\cdot x + a^{-1}\cdot(-a^2 x),\; a\cdot a^{-1}\bigr)
= (ax - ax,\; 1) = (0, 1).
\]

\end{proof}

It is at first surprising to come up with $\star$ operation, but $G$ is isomorphic to the direct product group. 
\begin{remark}[Group isomorphism to product group]
The map
\be
\Psi : G \to (X, +) \times (A, \cdot), \qquad
\Psi(x, a) := (ax, a),
\ee
is a group isomorphism onto the direct product of $(X, +)$ and $(A, \cdot)$ with operation
$\star_\Psi:(x,\; a) \star_\Psi (y,\; b) = (x+y,\; ab)$. Indeed,
\[
\Psi\bigl((x, a) \star (y, b)\bigr) = \bigl(ab\cdot(b^{-1}x + a^{-1}y),\; ab\bigr) = (ax + by,\; ab)
= \Psi(x, a) \star_\Psi \Psi(y, b),
\]
and $\Psi$ is bijective with inverse $\Psi^{-1}(y, b) = (b^{-1}y, b)$.

This isomorphism is used only to \emph{motivate} the characters below. The Lean development does not
formalize $\Psi$: it equips $G$ with its group structure (\lemref{lem:group}) and defines the characters
$\chi_{\alpha,k}$ (\propref{prop:group_characaters}) directly on the twisted group, proving the homomorphism
property by direct computation rather than by pulling back through $\Psi$. Accordingly it is recorded as a
remark, with no \texttt{\textbackslash lean} link and no node in the dependency graph.
\end{remark}

The isomorphism gives us a easy way to compute its irreducible representations which forms the Pontryagin dual group (the set of irrep characters for an abelian group).
\begin{remark}[Pontryagin dual over direct product]
For finite abelian groups $G_1, G_2$ with direct product $H = G_1 \times G_2$, the Pontryagin dual (the set of irreducible characters) factors as $\hat{H} = \hat{G_1} \times \hat{G_2}$. Like the isomorphism above, this is used only to motivate the character formula and is not part of the formalized graph.
\end{remark}
\begin{proposition}[Group characters]\label{prop:group_characaters}\lean{BlrPcp.gChar, BlrPcp.gChar_mul, BlrPcp.card_gChar_index, BlrPcp.gChar_injective}\leanok
The group $(G, \star)$ has the Pontryagin dual character group 
\be
\hat{G} = \{\chi_{\alpha, k}\}_{\alpha \in X, k \in [0, \cdots, q-2]}, \quad \chi_{\alpha, k}(x,a) = \omega^{\Tr \inner{\alpha}{a x}} \psi_k(a),
\ee
where $\psi_k(a) = e^{\frac{i 2 \pi m k}{q-1}}$ and $a=g_0^m$ for a fixed multiplicative generator $g_0\in A$.
\end{proposition}
\begin{proof}\leanok
By the isomorphism $\Psi(x, a) = (ax, a)$ in the remark above, $(G, \star) \cong (X, +) \times (A, \cdot)$. First we consider the characters of the product group. Using the Pontryagin-dual factorization in the remark above, we have
\[
\widehat{(X, +) \times (A, \cdot)} = \hat{X} \times \hat{A}.
\]
The character groups of $X$ and $A$ are
\begin{align}
    \hat{X} &= \{\phi_\alpha(y) = \omega^{\Tr \inner{\alpha}{y}}\}_{\alpha\in X}, \\
     \hat{A} &= \{\psi_k(b) =e^{\frac{i 2 \pi m k}{q-1}}\}_{k \in [0, \cdots, q-2]},
\end{align}
where $b=g_0^m$ for a fixed multiplicative generator $g_0\in A$. 
Here the field trace $\Tr: \F_q \to \F_p$ is defined by $\Tr(x) = x + x^p + x^{p^2} + \cdots + x^{p^{s-1}}$ for $x \in \F_q$ (see Definition~1.1 in Chapter 1).%(see Definition~\ref{def:field-trace} in \secref{sec:blr-general-finite-fields}).
Therefore, $\widehat{(X, +) \times (A, \cdot)} = \{\Theta_{\alpha, k}(y, b) = \phi_\alpha(y) \psi_k(b)\}_{\alpha, k}$.
The character $\chi \in \hat{G}$ is given by the pull-back through the isomorphism
\be
\chi_{\alpha, k}(x, a) = \Theta_{\alpha, k}(\Psi(x, a)) = \Theta_{\alpha, k}(ax, a).
\ee
Therefore, we have $\chi_{\alpha, k}(x, a) = \phi_\alpha(ax) \psi_k(a) =\omega^{\Tr \inner{\alpha}{a x}} \psi_k(a)$.
\end{proof}
\subsection{Approximate representation and matrix lift of $f$}\label{sec:rep}

\begin{definition}[Matrix lift $F_f$]\label{def:Flift}\lean{BlrPcp.liftMatrix}\leanok
Define a unitary matrix valued lift $F_f : G \to U(\cc^{q-1})$ by the diagonal matrix
\begin{equation}\label{eq:Flift-def}
F_f(x, a) = D(a f(x)) := \sum_{c \in A} \omega^{\Tr\, c a f(x)}\ket{c}\bra{c},
\end{equation}
where $\{\ket{c}\}$ denotes the orthonormal basis vectors of a Hilbert space $\mathcal{H} \cong \cc^{q-1}$ and the inner product between operators defined by $\inner{A}{B}_\sigma=\Tr(A^* B \sigma)$ and choose $\sigma = \I/(q-1)$.
\end{definition}

% \begin{lemma}[$F_f$ is unitary-valued]\label{lem:Ff-unitary}\uses{def:Flift}\lean{BlrPcp.liftMatrix_mem_unitaryGroup}\leanok
% For every $(x, a) \in G$, the lift $F_f(x, a)$ is a unitary matrix, i.e.\ $F_f : G \to U(\cc^{q-1})$.
% \end{lemma}
We first show that the BLR condition indeed gives an approximate represenatation.

\begin{lemma}[Lift operation from BLR condition]\label{lem:rep-iff-blr}\uses{lem:group, def:Flift}\lean{BlrPcp.liftFin_isApproxRepresentation}\leanok
Let $f : X \to \F_q$ and $(x, a), (y, b) \in G$ be rejected by the BLR test with probability $\varepsilon$,
then the unitary matrix lift $F_f$ in Equation~\ref{eq:Flift-def} is an $(\frac{q}{q-1}\varepsilon, \frac{\I}{q-1})$ approximate representation of $(G, \star)$.
\end{lemma}

\begin{proof}
The BLR test condition gives
\[
\Pr_{c,d,x,y}\!\bigl[c f(x) + d f(y) \neq f(cx + dy)\bigr] \leq \varepsilon.
\]
First we observe that
\begin{equation}\label{eq:rep-eq-blr}
\Pr_{g, h \in G}\!\bigl[F_f(g \star h) = F_f(g)F_f(h)\bigr]
= \Pr_{c,d,x,y}\!\bigl[c f(x) + d f(y) = f(cx + dy)\bigr].
\end{equation}
Using~Equation~\ref{eq:star} and~Equation~\ref{eq:Flift-def},
\begin{align*}
F_f\bigl((x,a)\star(y,b)\bigr) &= \sum_{c\in A}\omega^{\Tr\, cab\,f(b^{-1}x + a^{-1}y)} \ket{c}\bra{c}, \\
F_f(x, a)\,F_f(y, b) &= \sum_{c\in A} \omega^{\Tr\, c(af(x) + bf(y))} \ket{c}\bra{c}.
\end{align*}
These coincide iff their exponents agree modulo $p$ for $c \in A$: 
\[
ab\, f(b^{-1}x + a^{-1}y) = af(x) + bf(y) \mod p, \quad \forall c \in A.
\]
Therefore, we have equality over $\F_p$ 
\[
ab\, f(b^{-1}x + a^{-1}y) = af(x) + bf(y).
\]
Since $a, b$ have a multiplicative inverse, substitute $c := b^{-1}$, $d := a^{-1}$, we have
\[
f(cx + d y) = c f(x) + d f(y),
\]
which proves Equation~\ref{eq:rep-eq-blr}.

To relate the test condition to the approximate representation, first note the equivalent definition of approximate representation
\be
\mathbb{E}_{g, h \in G} \Re \left\langle F_f(g)^{*} F_f(h), F_f\left(g^{-1} h\right)\right\rangle_{\sigma} = \mathbb{E}_{g, h \in G} \Re\left\langle F_f(g) F_f(h), F_f\left(g h\right)\right\rangle_{\sigma},
\ee
because $F_f$ is unitary so $F_f(g)^* = F_f(g^{-1}), \forall g \in G$. Then evaluate the right hand side
\begin{align}
    \mathbb{E}_{g, h \in G} \Re \left\langle F_f(g) F_f(h), F_f\left(g h\right)\right\rangle_{\sigma} &= \mathbb{E}_{g, h \in G} \Re \Tr\left[F_f^*(h) F_f^*(g) F_f\left(g h\right)\sigma\right], \nonumber \\
    & = \mathbb{E}_{(x,a), (y,b) \in G} \frac{1}{q-1}\Re \Tr\left[ \sum_{c\in A} \omega^{-\Tr\, c(af(x) + bf(y))} \omega^{\Tr\, cab\,f(b^{-1}x + a^{-1}y)} \ket{c}\bra{c}\right], \nonumber \\
    & = \mathbb{E}_{(x,a), (y,b) \in G} \frac{1}{q-1}\Re \sum_{c\in A} \omega^{\Tr\, c z},
\end{align}
where we define $z=c (ab\,f(b^{-1}x + a^{-1}y) - (af(x) + bf(y)))$. To compute the expectation, we can split to two cases
\begin{align}
\sum_{c\in A} \omega^{\Tr\, c z} = 
    \begin{cases}
        q-1, \quad z=0 \quad \text{when $f$ passes BLR}, \\
        -1, \quad z \neq 0 \quad \text{when $f$ fails BLR}.
    \end{cases}
\end{align}
The second case is because $\sum_{c\in A} \omega^{\Tr\, c z} + 1= \sum_{c\in \FF_q} \omega^{\Tr\, c z} =0$. Therefore, we have the approximate representation
\be
\mathbb{E}_{g, h \in G} \Re \left\langle F_f(g) F_f(h), F_f\left(g h\right)\right\rangle_{\sigma} \geq (1- \varepsilon) - \frac{1}{q-1} \varepsilon = 1 - \frac{q}{q-1}\varepsilon. 
\ee
Hence $F_f$ is an approximate representation with  $\epsilon= \frac{q}{q-1}\varepsilon$ as claimed.

\end{proof}
The approximate representation ``appears'' slightly worse than what we would have hoped $\epsilon > \varepsilon$, but we will see this is not an issue for achieving a high soundness. Remarkably, the twisted group operation let us recover the linearity test condition exactly.

% \begin{fact}
%     $F_f$ is an exact representation of $(G, \star)$ if and only if
% $f \in \lin$.
% \end{fact}
% \begin{proof}
%     If $F_f$ is an exact representation of $G$, then~Equation~\ref{eq:rep-eq-blr} forces
% $\varepsilon = 0$. Setting $c = d = 1$ in the BLR condition gives $f(x + y) = f(x) + f(y)$
% for all $x, y$, so $f$ is additive. Set $y = 0$ and $c = 1 \in A$: $f(x) = f(x) + d\, f(0)$ for any
% $d \in A$. Set $y = 0$ and $d=1$, then $f(0)=0$ and $f(c x) = c f(x)$. Combined with additivity, $f$ is $\F_q$-linear, so $f \in \lin$.

% Conversely, if $f = \ell_\alpha$ for some $\alpha \in X$, then so $F_f$ is a group homomorphism and an exact representation of $G$.
% \end{proof}

\subsection{Gowers-Hatami implies BLR high soundness}
In this section, we first prove the key lemma relating the implication of GH theorem to the hamming distance. 
The proof outlined here correspond to the formalized version in lean. The most convenient proof still uses properties of the Fourier transform of $F_f$. However, we remark that there are other alternative proofs, which we leave as comment in the current source file.

% Finally, we provide a proof of \thmref{thm:main}.

% We provide two alternative proofs of the key lemma:
% \begin{enumerate}
%     \item The first proof in \lemref{lem:gh_linearity_1} should only use the existence of an isometry $V$ and exact representation $R$. Unfortunately, we ended up with an extra constraint about the range of $\epsilon$. 
%     \item In the second proof in \lemref{lem:gh_linearity_2}, we use an explicit isometry $V$ constructed via $F_f$ in the proof of Gowers-Hatami theorem. This proof might require a extra care in formalization of lean. However, this hybrid Fourier-GH approach gives us an unconditional high soundness of the linear BLR test.   
% \end{enumerate}

% \begin{lemma}[Gowers-Hatami to linearity (Convex character proof)] \label{lem:gh_linearity_1}
% Suppose $F_f: G \to U(\cc^{q-1})$ defined in Equation~\ref{eq:Flift-def} via $f : \FF_q^n \to \FF_q$ is an $(\epsilon, \sigma)$ approximate unitary representation of $(G, \star)$ defined in Equation~\ref{eq:star}. 
% % Furthermore, suppose $1 - \epsilon > \frac{\sqrt{q}}{q-1}$.
% Then, there exists a linear function $\ell: \FF_q^n \to \FF_q$ such that the relative Hamming distance $\delta(f, \ell)$ is bounded by:
% \be\label{eq:GH_linear_bound}
%     \delta(f, \ell) \le \frac{q-1}{q} \epsilon.
% \ee
% \end{lemma}

% \begin{proof}
% % \yan{In this specific proof below without any Fourier analysis, we also need the condition $1 - \epsilon > \frac{\sqrt{q}}{q-1}$, hence the additional assumption.}

%     By Gowers-Hatami theorem, we have an isometry $V: \mathcal{H} \to \mathcal{H}^\prime$ and an exact representation $R: G \to \End(\mathcal{H}^\prime)$, where the Hilbert space $\mathcal{H}\cong \cc^{q-1}$ and the space $R$ acts on $\mathcal{H}^\prime = L(G, \mathcal{H})$ 
%     \begin{equation}
%     \EE_{g \in G} \Re \langle F_f(g), V^* R(g) V \rangle_\sigma \ge 1 - \epsilon. \label{eq:inner1}
% \end{equation}
%     Let $\{ |c\rangle \}_{c \in \FF_q^\times}$ be the orthonormal basis for $\mathcal{H}$, and the set of characters $\{\chi\}_{\chi\in\hat{G}}$ be an orthonormal basis for $L(G, \cc)$. 
    
%     Because $F_f(g)$ is diagonal in $\{\ket{c}\}$ with entries $\omega^{\Tr\, caf(x)}$, we can rewrite the inner product trace:
%     \[
%         \langle F_f(g), V^* R(g) V \rangle_\sigma = \frac{1}{q-1} \sum_{c \in \FF_q^\times} \langle c | F_f^*(g) V^* R(g) V | c \rangle = \frac{1}{q-1} \sum_{c \in \FF_q^\times} \omega^{-\Tr\, caf(x)} \langle v_c, R(g) v_c \rangle_{L(G, \mathcal{H})},
%     \]
%     where $\ket{v_c} = V\ket{c} \in \mathcal{H}^\prime$.
%     Taking the expectation over $g = (x, a)$ and taking the real part, condition Equation~\ref{eq:inner1} becomes:
%     \begin{equation}
%         \frac{1}{q-1} \sum_{c \in \FF_q^\times} \EE_{x, a} \Re \left[ \omega^{-\Tr\, caf(x)} \langle v_c, R(x, a) v_c \rangle \right] \ge 1 - \epsilon. \label{eq:expanded}
%     \end{equation}

%     A basis for $L(G, \mathcal{H})$ is given by $\{\chi \ket{c}\}_{\chi\in\hat{G}, c\in \FF_q^\times}$. By \propref{prop:group_characaters}, the characters are given by
%     \[
%     \hat{G} = \{\chi_{\alpha, k}\}_{\alpha \in X, k \in [0, \cdots, q-2]}, \chi_{\alpha, k}(x,a) = \omega^{\Tr \inner{\alpha}{a x}} \psi_k(a).
%     \]
%     The regular representation $R(g)$ acts on $L(G, \mathcal{H})$ so it can be decomposed as 
%     \be\label{eq:reg_repr_decomp}
%     R(g) = \sum_{\alpha, k} \chi_{\alpha, k}(g) P_{\alpha, k},
%     \ee
%     where $P_{\alpha, k} = \ket{\chi_{\alpha, k}}\bra{\chi_{\alpha, k}} \otimes \I_{\mathcal{H}}$ projects onto the character function subspace (this leaves $\mathcal{H}$ invariant). 
    
%     Substituting this into the matrix element for $v_c$, we obtain:
%     \[
%         \langle v_c, R(g) v_c \rangle = \sum_{\alpha, k} \chi_{\alpha, k}(g) \langle v_c, P_{\alpha, k} v_c \rangle.
%     \]
%     Let $p_{c, \alpha, k} := \langle v_c, P_{\alpha, k} v_c \rangle$. As $P_{\alpha, k}$ is a projector, then 
%     \be
%     p_{c, \alpha, k} = \langle v_c \ket{\chi_{\alpha, k}}\bra{\chi_{\alpha, k}} \otimes \I_{\mathcal{H}}\ket{v_c} \geq 0.
%     \ee
    
%     Since $V$ is an isometry $V^* V = \I_{\mathcal{H}}$, we have $\bra{c} V^* V \ket{c} = \inner{c}{c}_{\mathcal{H}} = 1$. 
%     Equivalently, we have $\inner{v_c}{v_c}_{L(G, \mathcal{H})} = 1$. 
%     Since $\sum P_{\alpha, k} = I_{L(G,H)}$, it follows that $p_{c, \alpha, k}$ is a valid probability distribution over the characters,
%     \be
%     \sum_{\alpha, k} p_{c, \alpha, k} = 1, \quad \forall c \in \FF_q^\times.
%     \ee

%     Substituting the spectral decomposition Equation~\ref{eq:reg_repr_decomp} back into Equation~\ref{eq:expanded}, we get:
%     \[
%         \sum_{c \in \FF_q^\times} \left( \frac{1}{q-1} \right) \sum_{\alpha, k} p_{c, \alpha, k} \Re \left( \EE_{x, a} \left[ \omega^{-\Tr\, caf(x)} \chi_{\alpha, k}(x, a) \right] \right) \ge 1 - \epsilon.
%     \]
%     This is a convex combination (a weighted average) of real values less than $1$. Because the average is at least $1 - \epsilon$, there must exist at least one specific configuration $(c^*, \alpha^*, k^*)$ with $c^* \neq 0$ such that its individual value satisfies:
%     \begin{equation}
%         \Re \left( \EE_{x, a} \left[ \omega^{-\Tr\, c^*af(x)} \omega^{\Tr \langle \alpha^*, ax \rangle} \psi_{k^*}(a) \right] \right) \ge 1 - \epsilon. \label{eq:dominant}
%     \end{equation}

%     Because $c^* \in \FF_q^\times$ is invertible, we define a specific linear function $\ell(x) := \langle (c^*)^{-1} \alpha^*, x \rangle$. Notice that $\langle \alpha^*, ax \rangle = c^* a \ell(x)$. We can rewrite the argument of the expectation in Equation~\ref{eq:dominant} as:
%     \[
%         E^* := \EE_{x} \left[ \EE_{a} \left[ \omega^{\Tr\, c^* a (\ell(x) - f(x))} \psi_{k^*}(a) \right] \right].
%     \]

%     For the rest of this proof, without using specific form of $p_{c, \alpha, k}$, we consider two cases: $k^*=0$ and $k^* \neq 0$. Define $z(x) = c^*(\ell(x) - f(x))$.

%     \begin{enumerate}
%         \item Consider $k^*=0$ so that $\psi_0(a) = 1$, the inner expectation over $a \in \FF_q^\times$ for a fixed $x$:
%     \be
%     \EE_{a \in \in \FF_q^\times}\omega^{\Tr\, c^* a (\ell(x) - f(x))} = 
%     \begin{cases}
%         1, \quad z(x) = 0, \\
%         - \frac{1}{q-1}, \quad z(x) \neq 0,
%     \end{cases}
%     \ee
%     where in the second case we have again used average of roots of unity is $\sum_{b \in \in \FF_q^\times} \omega^{\Tr b} +1 = 0$.

%     Evaluating the outer expectation over $x$ and recall the definition of hamming distance 
%     \be
%     \delta_\alpha := \Pr_x[f(x) \ne \ell_\alpha(x)],
%     \ee
%     we find:
%     \[
%         \Re E^* = E^* = 1 \cdot (1 - \delta(f, \ell)) + \left( \frac{-1}{q-1} \right) \cdot \delta(f, \ell) = 1 - \delta(f, \ell) \left( 1 + \frac{1}{q-1} \right) = 1 - \delta(f, \ell) \left( \frac{q}{q-1} \right).
%     \]
%     Applying our bound from Equation~\ref{eq:dominant}, we have:
%     \[
%         1 - \delta(f, \ell) \left( \frac{q}{q-1} \right) \ge 1 - \epsilon.
%     \]
%     Rearranging yields the final result:
%     \[
%         \delta(f, \ell) \le \frac{q-1}{q} \epsilon.
%     \]

%     \item Consider $k^* \neq 0$, then by \propref{prop:group_characaters} we have $\psi_k(a) = e^{\frac{i 2 \pi m k}{q-1}}$ and $a=g_0^m$ for a fixed multiplicative generator $g_0\in \FF_q^\times$. Note that averaging over $\psi_k(a)$ we have
%     \be\label{eq: averaging_psi}
%     \EE_{a \in \FF_q^\times} \psi_k(a) = \frac{1}{q-1} \sum_{m=0}^{q-2} e^{\frac{i 2 \pi m k}{q-1}} = 0.
%     \ee
%     The average expectation over $a \in \FF_q^\times$ for a fixed $x$ gives:
%     \be
%     \EE_{a \in \in \FF_q^\times}\omega^{\Tr \, c^* a (\ell(x) - f(x))}  \psi_{k^*}(a) = 
%     \begin{cases}
%         0, \quad z(x) = 0, \\
%         \frac{1}{q-1} \sum_{a \in \in \FF_q^\times} \omega^{\Tr a z(x)}\psi_{k^*}(a), \quad z(x) \neq 0.
%     \end{cases}
%     \ee
%     In the second case, let $b = a z(x)$ so that
%     \begin{align}
%         \frac{1}{q-1} \sum_{a \in \in \FF_q^\times} \omega^{\Tr\, a z(x)}\psi_{k^*}(a) &= \frac{1}{q-1} \sum_{b \in \in \FF_q^\times} \omega^{\Tr b}\psi_{k^*}(z^{-1}b), \nonumber \\
%         & =  \frac{\psi_{k^*}(z^{-1})}{q-1} \sum_{b \in \in \FF_q^\times} \omega^{\Tr b}\psi_{k^*}(b).
%     \end{align}
%     It turns out the summation is a known quantitiy called \textit{the Gauss sum}
%     \be
%     \mathcal{G}(\psi_{k^*}) = \sum_{b \in \in \FF_q^\times} \omega^{\Tr b}\psi_{k^*}(b).
%     \ee
%     Is it believed that there is no analytic expression for $\mathcal{G}$ but it is known that
%     \be\label{eq:gauss_sum_mag}
%     \abs{\mathcal{G}(\psi_{k^*})} = \sqrt{q}
%     \ee

%     Evaluating the outer expectation over $x$ and recall the definition of hamming distance gives
%     \[
%     \Re E^* = 0 \cdot (1 - \delta(f, \ell)) + \Re \left( \frac{\psi_{k^*}(z^{-1})}{q-1} \mathcal{G}(\psi_{k^*})\right) \cdot \delta(f, \ell) .
%     \]
%     Therefore we have
%     \[
%     \Re E^* = \Re \left( \frac{\psi_{k^*}(z^{-1})}{q-1} \mathcal{G}(\psi_{k^*})\right) \delta(f, \ell) \leq \frac{\abs{\mathcal{G}(\psi_{k^*})}}{q-1} = \frac{\sqrt{q}}{q-1}.
%     \]
%     On the other hand, from Equation~\ref{eq:dominant} we have
%     \[
%     \Re E^* \geq 1- \epsilon.
%     \]

%     This is a contradiction to the existence of $k^*$ if 
%     \be\label{eq:epsilon_bound}
%     \frac{\sqrt{q}}{q-1} < 1- \epsilon.
%     \ee
%     Then if we have a good test outcome such that Equation~\ref{eq:epsilon_bound} is satisfied, then we only have a contribution from $k^*=0$, where we have already achieved the bound.
%     \end{enumerate}
    
% \end{proof}

% The above bound is tight assuming we pass the test with high probability such that $\epsilon < 1 - \frac{\sqrt{q}}{q-1} $. Next, we try to give a proof that is unconditional, using Fourier coefficients. 

First let us write down statements about Fourier coefficients of $F_f$, which we will use in the proof.

\begin{proposition}[Fourier coefficient of $F_f$]\label{prop:Ff_fourier}\uses{def:Flift, prop:group_characaters}\lean{BlrPcp.gFourierCoeff_factor}\leanok
    Let the matrix lift $F_f : G \to U(\cc^{q-1})$ be defined by the diagonal matrix
\begin{equation*}
F_f(x, a) = D(a f(x)) := \sum_{c \in A} \omega^{\Tr\, c a f(x)}\ket{c}\bra{c},
\end{equation*}
then its Fourier coefficients can be written as
\be\label{eq:factored_coeff}
[\hat{F}_f(\chi_{c\beta, k})]_{cc} = \psi_k(c) \nu_k(\beta),
\ee
where 
\be\label{def:nu_k}
\nu_k(\beta) := \EE_{x,b} \left[ \omega^{\Tr b(f(x) - \langle \beta, x \rangle)} \overline{\psi_k(b)} \right],
\ee

\end{proposition}

\begin{proof}
    By \propref{prop:group_characaters}, the characters are given by
    \[
    \hat{G} = \{\chi_{\alpha, k}\}_{\alpha \in X, k \in [0, \cdots, q-2]}, \chi_{\alpha, k}(x,a) = \omega^{\Tr \inner{\alpha}{a x}} \psi_k(a).
    \]
    Therefore the Fourier coefficients of $F_f$ are
    \begin{align}
        \hat{F}_f(\chi_{\alpha, k}) 
        & = \sum_{c \in \FF_q^\times } \EE_{x, a} \omega^{\Tr\, c a f(x)} \bar{\chi}_{\alpha, k}(x, a) , \nonumber \\
        & = \sum_{c \in \FF_q^\times }\EE_{x,a} \left[ \omega^{\Tr\, a(cf(x) - \langle \alpha, x \rangle)} \overline{\psi_k(a)} \right] \ket{c}\bra{c}.
    \end{align}
    

Since $c \in \FF_q^\times$ is invertible, we perform a change of variables. Let $\alpha = c\beta$ for a fresh $\beta \in \FF_q^n$. Because $c$ is fixed, this is a bijection on $\FF_q^n$. 
We also substitute $b = ca$, which is a bijection on $\FF_q^\times$. 
This means $a = bc^{-1}$, and by the multiplicity of Dirichlet characters, $\overline{\psi_k(a)} = \overline{\psi_k(bc^{-1})} = \overline{\psi_k(b)}\psi_k(c)$.

Substituting these into the expectation isolates the $c$-dependence: 
\begin{align} 
    [\hat{F}_f(\chi_{c\beta, k})]_{cc} &= \EE_{x,b} \left[ \omega^{\Tr b(f(x) - \langle \beta, x \rangle)} \overline{\psi_k(b)}\psi_k(c) \right] \nonumber \\
    &= \psi_k(c) \nu_k(\beta), 
\end{align}
where $\nu_k(\beta) := \EE_{x,b} \left[ \omega^{\Tr b(f(x) - \langle \beta, x \rangle)} \overline{\psi_k(b)} \right]$ is a $c$-independent coefficient.    
\end{proof}


% \begin{lemma}[Gowers-Hatami to linearity (Convex character Fourier hybrid)] \label{lem:gh_linearity_2}
% Suppose $F_f: G \to U(\cc^{q-1})$ defined in Equation~\ref{eq:Flift-def} via $f : \FF_q^n \to \FF_q$ is an $(\epsilon, \sigma)$ approximate unitary representation of $(G, \star)$ defined in Equation~\ref{eq:star}. 
% Then, there exists a linear function $\ell: \FF_q^n \to \FF_q$ such that the relative Hamming distance $\delta(f, \ell)$ is bounded by:
% \be\label{eq:GH_linear_bound}
%     \delta(f, \ell) \le \frac{q-1}{q} \epsilon.
% \ee
% \end{lemma}

% \begin{proof}

%     By Gowers-Hatami theorem, we have an isometry $V: \mathcal{H} \to \mathcal{H}^\prime$ and an exact representation $R: G \to \End(\mathcal{H}^\prime)$, where the Hilbert space $\mathcal{H}\cong \cc^{q-1}$ and the space $R$ acts on $\mathcal{H}^\prime = L(G, \mathcal{H})$
%     \begin{equation}
%     \EE_{g \in G} \Re \langle F_f(g), V^* R(g) V \rangle_\sigma \ge 1 - \epsilon. \label{eq:inner1}
% \end{equation}
%     Let $\{ |c\rangle \}_{c \in \FF_q^\times}$ be the orthonormal basis for $\mathcal{H}$, and the set of characters $\{\chi\}_{\chi\in\hat{G}}$ be an orthonormal basis for $L(G, \cc)$. 
    
%     % Since $V$ is an isometry, the vectors $\{v_c := V|c\rangle\}$ form an orthonormal set of $q-1$ basis in $L(G, \mathcal{H})$. 
    
%     Because $F_f(g)$ is diagonal in $\{\ket{c}\}$ with entries $\omega^{\Tr\, caf(x)}$, we can rewrite the inner product trace:
%     \[
%         \langle F_f(g), V^* R(g) V \rangle_\sigma = \frac{1}{q-1} \sum_{c \in \FF_q^\times} \langle c | F_f^*(g) V^* R(g) V | c \rangle = \frac{1}{q-1} \sum_{c \in \FF_q^\times} \omega^{-\Tr\, caf(x)} \langle v_c, R(g) v_c \rangle_{L(G, \mathcal{H})},
%     \]
%     where $\ket{v_c} = V\ket{c} \in \mathcal{H}^\prime$.
%     Taking the expectation over $g = (x, a)$ and taking the real part, condition Equation~\ref{eq:inner1} becomes:
%     \begin{equation}
%         \frac{1}{q-1} \sum_{c \in \FF_q^\times} \EE_{x, a} \Re \left[ \omega^{-\Tr\, caf(x)} \langle v_c, R(x, a) v_c \rangle \right] \ge 1 - \epsilon. \label{eq:expanded}
%     \end{equation}

%     A basis for $L(G, \mathcal{H})$ is given by $\{\chi \ket{c}\}_{\chi\in\hat{G}, c\in \FF_q^\times}$. By \propref{prop:group_characaters}, the characters are given by
%     \[
%     \hat{G} = \{\chi_{\alpha, k}\}_{\alpha \in X, k \in [0, \cdots, q-2]}, \chi_{\alpha, k}(x,a) = \omega^{\Tr \inner{\alpha}{a x}} \psi_k(a).
%     \]

%     Moreover, the Fourier coefficients of $F_f$ are
%     \begin{align}
%         \hat{F}_f(\chi_{\alpha, k}) 
%         & = \sum_{c \in \FF_q^\times } \EE_{x, a} \omega^{\Tr\, caf(x)} \bar{\chi}_{\alpha, k}(x, a) , \nonumber \\
%         & = \sum_{c \in \FF_q^\times }\EE_{x,a} \left[ \omega^{\Tr\, a(cf(x) - \langle \alpha, x \rangle)} \overline{\psi_k(a)} \right] \ket{c}\bra{c}.
%     \end{align}
    
%     The regular representation $R(g)$ acts on $L(G, \mathcal{H})$ so it can be decomposed as 
%     \be\label{eq:reg_repr_decomp}
%     R(g) = \sum_{\alpha, k} \chi_{\alpha, k}(g) P_{\alpha, k},
%     \ee
%     where $P_{\alpha, k} = \ket{\chi_{\alpha, k}}\bra{\chi_{\alpha, k}} \otimes \I_{\mathcal{H}}$ projects onto the character function subspace (this leaves $\mathcal{H}$ invariant). 
    
%     Substituting this into the matrix element for $v_c$, we obtain:
%     \[
%         \langle v_c, R(g) v_c \rangle = \sum_{\alpha, k} \chi_{\alpha, k}(g) \langle v_c, P_{\alpha, k} v_c \rangle.
%     \]
%     Let $p_{c, \alpha, k} := \langle v_c, P_{\alpha, k} v_c \rangle$. As $P_{\alpha, k}$ is a projector, then 
%     \be
%     p_{c, \alpha, k} = \langle v_c \ket{\chi_{\alpha, k}}\bra{\chi_{\alpha, k}} \otimes \I_{\mathcal{H}}\ket{v_c} \geq 0.
%     \ee
%     Let the isometry be given by $V\ket{c} = [F_f(g)]_{c}$, then the matrix element is given by the Fourier coefficients
%     \be
%     p_{c, \alpha,k} = \abs{[\hat{F}_f(\chi_{\alpha, k})]_{cc}}^2.
%     \ee
%     By \propref{prop:Ff_fourier}, let $\alpha = c \beta$, we can show
%     \be
%     p_{c, c\beta,k} = \abs{\psi_k(c) \nu_k(\beta)}^2 = \abs{\nu_k(\beta)}^2,
%     \ee
%     where 
%     \[
%     \nu_k(\beta) := \EE_{x,b} \left[ \omega^{\Tr\, b(f(x) - \langle \beta, x \rangle)} \overline{\psi_k(b)} \right],
%     \]
%     is defined in Equation~\ref{def:nu_k}.   
    
%     Since $V$ is an isometry $V^* V = \I_{\mathcal{H}}$, we have $\bra{c} V^* V \ket{c} = \inner{c}{c}_{\mathcal{H}} = 1$. 
%     Equivalently, we have $\inner{v_c}{v_c}_{L(G, \mathcal{H})} = 1$. 
%     Since $\sum P_{\alpha, k} = I_{L(G,H)}$, it follows that $p_{c, \alpha, k}$ is a valid probability distribution over the characters,
%     \be
%     \sum_{\alpha, k} p_{c, \alpha, k} = 1, \quad \forall c \in \FF_q^\times.
%     \ee

%     Equivalently, we have the relabeled probability
%     \be
%     \sum_{\beta, k} p_{c, c \beta, k} = 1, \quad \forall c \in \FF_q^\times.
%     \ee
%     Therefore $|\nu_k(\beta)|^2$ is a probability distribution indexed by $(k, \beta)$. Because the $k=0$ probability is a subset, we have
%     \be\label{eq:nu_less_1}
%     \sum_{\beta \in \FF_q^n} |\nu_0(\beta)|^2 \leq \sum_{\beta \in \FF_q^n} \sum_k |\nu_k(\beta)|^2 = 1.
%     \ee
    
%     Substituting the spectral decomposition Equation~\ref{eq:reg_repr_decomp} back into Equation~\ref{eq:expanded}, we define:
%     \begin{align*}
%             E &:= \frac{1}{q-1} \sum_{c \in \FF_q^\times} \EE_{x, a} \Re \left[ \omega^{-\Tr\, caf(x)} \langle v_c, R(x, a) v_c \rangle \right], \\
%             % & = \sum_{c \in \FF_q^\times} \left( \frac{1}{q-1} \right) \sum_{\alpha, k} p_{c, \alpha, k} \Re \left( \EE_{x, a} \left[ \omega^{-\Tr\, caf(x)} \chi_{\alpha, k}(x, a) \right] \right), \\
%             & = \frac{1}{q-1} \sum_{c \in \FF_q^\times}\sum_{\alpha, k} p_{c, \alpha, k} \Re \left( \EE_{x, a} \left[ \omega^{-\Tr\, caf(x)} \omega^{\Tr\, \inner{\alpha}{ax}} \psi_k(a) \right] \right).
%     \end{align*}

%     Relabel the sum by $\alpha = c \beta$ and $b=ca$, we have
%     \[
%     E =  \frac{1}{q-1} \sum_{c \in \FF_q^\times}\sum_{\beta, k} \abs{\nu_k(\beta)}^2 \Re \left( \EE_{x, b} \left[ \omega^{\Tr\, b(\inner{\beta}{x} - f(x))} \psi_k(c^{-1})\psi_k(b) \right] \right).
%     \]
%     Use the indicator function of $\FF_q^\times$:
%     \[
%         \sum_{c \in \FF_q^\times} {\psi_k(c^{-1})} = (q-1) \I(k=0)
%     \] 
%     we have     
%     \[
%     E = \sum_{\beta} \abs{\nu_0(\beta)}^2 \Re \left( \EE_{x, b} \left[ \omega^{\Tr\, b(\inner{\beta}{x} - f(x))} \right] \right).
%     \]
    
%     This is a convex combination (a weighted average) of real values less than $1$ with the weight sum to less than 1 by Equation~\ref{eq:nu_less_1}. Because the average is at least $E \geq 1 - \epsilon$, there must exist at least one specific configuration $\beta^*$ such that its individual value satisfies:
%     \begin{equation}
%         E^* := \Re \left( \EE_{x, b} \left[ \omega^{\Tr\, b(\inner{\beta^*}{x} - f(x))} \right] \right) \ge 1 - \epsilon. \label{eq:dominant}
%     \end{equation}
    
%     Let $\ell_\beta(x):=\inner{\beta^*}{x}$, the inner expectation over $b \in \FF_q^\times$ for a fixed $x$:
%         \be
%         \EE_{b \in \in \FF_q^\times}\omega^{\Tr\, b z(x)} = 
%         \begin{cases}
%             1, \quad z(x) = 0, \\
%             - \frac{1}{q-1}, \quad z(x) \neq 0,
%         \end{cases}
%         \ee
%         where in the second case we have again used average of roots of unity is $\sum_{b \in \in \FF_q^\times} \omega^{\Tr\, b} +1 = 0$.
    
%         Evaluating the outer expectation over $x$ and recall the definition of hamming distance 
%         \be
%         \delta(f, \ell_\alpha) := \Pr_x[f(x) \ne \ell_\alpha(x)],
%         \ee
%         we find:
%         \[
%             E^* = 1 \cdot (1 - \delta(f, \ell_\beta)) + \left( \frac{-1}{q-1} \right) \cdot \delta(f, \ell_\beta) = 1 - \delta(f, \ell_\beta) \left( 1 + \frac{1}{q-1} \right) = 1 - \delta(f, \ell_\beta) \left( \frac{q}{q-1} \right).
%         \]
%         Applying our bound from Equation~\ref{eq:dominant}, we have:
%         \[
%             1 - \delta(f, \ell_\beta) \left( \frac{q}{q-1} \right) \ge 1 - \epsilon.
%         \]
%         Rearranging yields the final result: $\exists \ell_\beta = \inner{\beta^*}{x}$ such that 
%         \[
%             \delta(f, \ell_\beta) \le \frac{q-1}{q} \epsilon.
%         \]
    
% \end{proof}


\begin{lemma}[Gowers-Hatami to linearity (formalized proof)] \label{lem:gh_linearity_formalized}\uses{prop:group_characaters, prop:Ff_fourier, gowers-hatami-correlation}\lean{BlrPcp.gh_linearity}\leanok
Suppose $F_f: G \to U(\cc^{q-1})$ defined in Equation~\ref{eq:Flift-def} via $f : \FF_q^n \to \FF_q$ is an $(\epsilon, \sigma)$ approximate unitary representation of $(G, \star)$ defined in Equation~\ref{eq:star}. 
Implicitly, $\epsilon \leq 1$.

Then, there exists a linear function $\ell: \FF_q^n \to \FF_q$ such that the relative Hamming distance $\delta(f, \ell)$ is bounded by:

\be\label{eq:GH_linear_bound}
\delta(f, \ell) \le \frac{q-1}{q} \epsilon.
\ee
\end{lemma}

\begin{proof}
Write the diagonal lift as $F_f(g) = \operatorname{diag}_{c \in \FF_q^\times}(\varphi_c(g))$, where the scalar functions are $\varphi_c(x,a) = \omega^{\Tr(caf(x))}$. 

Because $F_f$ is diagonal and $\sigma = \I/(q-1)$, the trace equates to the average over the diagonal elements. 
By an intermediate statement in the proof of Gowers-Hatami theorem using Lemma.~\ref{gowers-hatami-correlation}, the compression from an exact representation can be rewritten as a pair-correlation:
\begin{equation}
    \mathcal{C}(f) := \EE_{g,h} \Re \inner{F_f(g) F_f(h)}{F_f(gh)}_\sigma \ge 1 - \epsilon. \label{eq:gh_pair_correlation}
\end{equation}


To connect the operator pair-correlation directly to the cubic trace of the Fourier coefficients, we express the lift matrices using the inverse Fourier transform, $F_f(x) = \sum_{\chi \in \widehat{G}} \hat{F}_f(\chi) \chi(x)$. Substituting this into the inner product definition of the pair-correlation and using the character homomorphism $\chi(gh) = \chi(g)\chi(h)$, we expand the trace:
\begin{align}
    \mathcal{C}(f) &= \EE_{g,h \in G} \Re \Tr\left( F_f(g)^* F_f(h)^* F_f(gh) \sigma \right) \nonumber \\
    &= \Re \Tr \left( \sum_{\chi_1, \chi_2, \chi_3} \hat{F}_f^*(\chi_1) \hat{F}_f^*(\chi_2) \hat{F}_f(\chi_3) \sigma \cdot \EE_g\left[\overline{\chi_1(g)}\chi_3(g)\right] \EE_h\left[\overline{\chi_2(h)}\chi_3(h)\right] \right).
\end{align}
By the orthogonality of characters, the expectations over $g$ and $h$ independently vanish unless $\chi_1 = \chi_3$ and $\chi_2 = \chi_3$. This completely collapses the triple sum to a single sum over $\chi$. Because the Fourier coefficients of the diagonal lift $F_f$ are themselves diagonal matrices (and thus commute), we can rearrange the surviving terms to recover the exact cubic trace expression:
\begin{equation}
    \mathcal{C}(f) = \Re \sum_{\chi \in \widehat{G}} \Tr\left( \hat{F}_f^*(\chi) \hat{F}_f(\chi) \hat{F}_f^*(\chi) \sigma \right).
\end{equation}

    Because the Fourier coefficients are all fully diagonal since $F_f$ is diagonal, we expand the product
    \[
    \sum_{\chi_{\alpha, k} \in \widehat{G}} \Tr(\hat{F}_f^*(\chi)\hat{F}_f(\chi)\hat{F}_f^*(\chi) \sigma) = \frac{1}{q-1}\sum_{\alpha \in \FF_q^n} \sum_{k=0}^{q-2} \sum_{c \in \FF_q^\times} \left| [\hat{F}_f(\chi_{\alpha, k})]_{cc} \right|^2 \overline{[\hat{F}_f(\chi_{\alpha, k})]_{cc}}.
    \]

    Using the substitution $\alpha = c\beta$, summing over all $\alpha$ is equivalent to summing over all $\beta$. 
Substituting our factored coefficient from Eq. \eqref{eq:factored_coeff} in \propref{prop:Ff_fourier}: 
    \begin{align*}
        \sum_{\chi \in \widehat{G}} \Tr(\hat{F}_f^*(\chi)\hat{F}_f(\chi)\hat{F}_f^*(\chi) \sigma) &= \frac{1}{q-1}\sum_{c \in \FF_q^\times} \sum_{\beta \in \FF_q^n} \sum_{k=0}^{q-2} |\psi_k(c)\nu_k(\beta)|^2 \overline{\psi_k(c)\nu_k(\beta)} \\
        &= \frac{1}{q-1}\sum_{\beta, k} |\nu_k(\beta)|^2 \overline{\nu_k(\beta)} \sum_{c \in \FF_q^\times} \overline{\psi_k(c)},
    \end{align*}
    where 
    \[
    \nu_k(\beta) := \EE_{x,b} \left[ \omega^{b(f(x) - \langle \beta, x \rangle)} \overline{\psi_k(b)} \right],
    \]
    is defined in \equref{def:nu_k}.

    Use the indicator function of $\FF_q^\times$:
    \[
        \sum_{c \in \FF_q^\times} \overline{\psi_k(c)} = (q-1) \I(k=0)
    \] 
we have
    \be\label{eq:triple_F_nu0}
        \sum_{\chi \in \widehat{G}} \Tr(\hat{F}_f^*(\chi)\hat{F}_f(\chi)\hat{F}_f^*(\chi) \sigma) = \sum_{\beta \in \FF_q^n} |\nu_0(\beta)|^2 \overline{\nu_0(\beta)}.
    \ee
    Note that $\nu_0(\beta) \in \mathbb{R}$ is related to the hamming distance $\delta(f, \ell_\beta)$ for liner function $\ell_\beta = \inner{f}{\ell_\beta}$ via
    \be\label{eq:nu0_hamming}
    \nu_0(\beta) = (1 - \delta(f, \ell_\beta)) - \frac{1}{q-1}\delta(f, \ell_\beta) = 1 - \frac{q}{q-1}\delta(f, \ell_\beta).
    \ee
    Moreover, because the Fourier coefficients sum to identity by \equref{eq:F_parseval}, we have
    \begin{align}
        1& = \frac{1}{q-1}\sum_{\chi\in\hat{G}}\Tr\left[\hat{F}^*_f(\chi)\hat{F}_f(\chi)\right] , \nonumber \\
        &= \frac{1}{q-1}\sum_{c \in \FF_q^\times} \sum_{\beta \in \FF_q^n} \sum_k |\psi_k(c)\nu_k(\beta)|^2, \nonumber \\
        &= \sum_{\beta \in \FF_q^n} \sum_k |\nu_k(\beta)|^2.
    \end{align}

    Therefore $|\nu_k(\beta)|^2$ is a probability distribution indexed by $(k, \beta)$. Because the $k=0$ probability is a subset, we have
    \be
    \sum_{\beta \in \FF_q^n} |\nu_0(\beta)|^2 \leq \sum_{\beta \in \FF_q^n} \sum_k |\nu_k(\beta)|^2 = 1.
    \ee

    Let $\beta^*$ be the closest linear function such that
    \be
    \nu_0(\beta^*) := \max_{\beta \in \FF_q^n} \nu_0(\beta).
    \ee
    Therefore from Gowers-Hatami implication in \equref{eq:gh_pair_correlation} and the simplification of Fourier coefficeints in \equref{eq:triple_F_nu0}, we have the bound
    \begin{align}
        1 - \epsilon \leq \Re \sum_{\chi\in\widehat G} \Tr(\hat{F}_f^*(\chi)\hat{F}_f(\chi)\hat{F}_f^*(\chi) \sigma) = \sum_{\beta \in \FF_q^n} |\nu_0(\beta)|^2 \overline{\nu_0(\beta)} \leq  \nu_0(\beta^*) \sum_{\beta \in \FF_q^n} |\nu_0(\beta)|^2 \leq  \nu_0(\beta^*).
    \end{align}

    Unpack the relation to the hamming distance, we have
    \be
    1 - \frac{q}{q-1} \delta(f, \ell_{\beta^*}) \ge 1 - \epsilon \implies \delta(f, \ell_{\beta^*}) \le \frac{q-1}{q}\epsilon.
    \ee

\end{proof}

% Finally, we are ready to prove the main theorem of BLR linearity test soundness. 

% \setcounter{theorem}{0}
% \begin{theorem}[Soundness of the BLR test]\label{thm:main}\uses{lem:rep-iff-blr, lem:gh_linearity_formalized}\lean{BlrPcp.blr_soundness}\leanok
% Let $q$ be prime. For every $f : \F_q^n \to \F_q$, $\delta(f, \lin) \le \Pr_{a,b,x,y}\!\bigl[af(x) + bf(y) \ne f(ax + by)\bigr]$.
% \end{theorem}    

% \begin{proof}
% First denote 
% \[
% X := \F_q^n, \qquad A := \F_q^\times, \qquad \omega := e^{2\pi i/p}.
% \]
% Define a group $G$ by the set $X \times A$ with the binary operation
% \begin{equation}
% (x, a) \star (y, b) := \bigl(b^{-1}x + a^{-1}y,\; ab\bigr).
% \end{equation}

%     Let $F_f$ be a unitary matrix valued lift $F_f : G \to U(\cc^{q-1})$ by the diagonal matrix
% \begin{equation}
% F_f(x, a) = D(a f(x)) := \sum_{c \in A} \omega^{\Tr\, c a f(x)}\ket{c}\bra{c},
% \end{equation}
% where $\{\ket{c}\}$ denotes the orthonormal basis vectors of a Hilbert space $\mathcal{H} \cong \cc^{q-1}$ and the inner product between operators defined by $\inner{A}{B}_\sigma=\Tr(A^* B \sigma)$ and choose $\sigma = \I/(q-1)$. By \lemref{lem:rep-iff-blr}, since $f$ is rejected with probability $\varepsilon$, then $F_f$ is an $(\frac{q}{q-1}\varepsilon,\frac{\I}{q-1})$ approximate representation of $(G, \star)$.

% Then by \lemref{lem:gh_linearity_3}, there exists a linear function $\ell: \FF_q^n \to \FF_q$ such that the relative Hamming distance $\delta(f, \ell)$ is bounded by:
% \be
%     \delta(f, \ell) \le \frac{q-1}{q} \frac{q}{q-1}\varepsilon = \varepsilon.
% \ee

% This proves the claim
% \[
% \delta(f, \lin) := \min_{\alpha} \delta(f, \ell_\alpha) \leq \delta(f, \ell) \leq \varepsilon.
% \]
% \end{proof}

% Miraculously, the loose factor in $(\epsilon^\prime =\frac{q}{q-1}\varepsilon,\frac{\I}{q-1})$ approximate representation $F_f$ cancels with the tighter factor in 
% \[
% \delta(f, \ell) \le \frac{q-1}{q} \epsilon^\prime.
% \]
